\section*{第 \ref{ch:b3:holomorphic} 章——全纯
函数}

\begin{solution}{exo:b3:holomorphic:1}
(a) $\bar z$：$P = x$，$Q = -y$，则 $P_x = 1 \neq -1 = Q_y$：
处处不 $\C$-可导。$\abs z^2$：$P = x^2 + y^2$，$Q =
0$：\omterm{prop:b3:holomorphic:cauchyriemann}{柯西–黎曼方程}要求 $2x = 0 = 2y$：可导性仅在
$0$ 处成立——从而全纯性处处不成立（不含任何\omterm{def:b3:topology:topology}{开集}）。
$\operatorname{Re}z$：$P_x = 1 \neq 0 = Q_y$：处处不可导。

(b) $x^2 - y^2 = \operatorname{Re}(z^2)$：$f(z) = z^2$ 满足要求。
所有解：若 $\operatorname{Re}f = \operatorname{Re}g$
其中 $f, g$ 在连通的 $\C$ 上全纯，则 $h = f -
g$ 满足 $\operatorname{Re}h = 0$；由\omterm{prop:b3:holomorphic:cauchyriemann}{柯西–黎曼方程}得 $h' =
P_x + \iu Q_x = 0 - \iu P_y = 0$：$h$ 为纯虚
常数。故 $f(z) = z^2 + \iu c$，$c \in \R$。
\end{solution}

\begin{solution}{exo:b3:holomorphic:2}
在单位圆 $\gamma(t) = \eu^{\iu t}$ 上：
\[
\int_C z^n\,\dd z = \int_0^{2\pi}\eu^{\iu nt}\,\iu\eu^{\iu
t}\dd t = \iu\int_0^{2\pi}\eu^{\iu(n+1)t}\dd t
= \begin{cases} 2\iu\pi & n = -1,\\ 0 & n \neq -1.\end{cases}
\]
对于 $\bar z$，沿 $[0, 1+\iu]$，$\gamma(t) = t(1 + \iu)$：
$\int_0^1 t(1 - \iu)(1 + \iu)\dd t = \int_0^12t\,\dd t = 1$。
沿 $0 \to 1 \to 1 + \iu$：$\int_0^1t\dd t +
\int_0^1(1 - \iu t)\,\iu\,\dd t = \frac12 + \iu + \frac12 = 1
+ \iu$。两条路径端点相同而积分值不同，由
\cref{prop:b3:holomorphic:primitive}，$\bar z$ 在任何包含两条路径的\omterm{def:b3:topology:topology}{开集}上
都没有原函数。
\end{solution}

\begin{solution}{exo:b3:holomorphic:3}
(a) 割去负实轴后的复平面 $\Omega = \C\setminus\intoc{-\infty}0$ 是
关于 $1$ 星形，且 $\frac1z \in \mathcal H(\Omega)$：
由\cref{thm:b3:holomorphic:cauchy}可知存在原函数 $L$
满足 $L(1) = 0$。于是 $\bigl(z\eu^{-L(z)}\bigr)' =
\eu^{-L}(1 - z\cdot\frac1z) = 0$：$z = c\,\eu^{L(z)}$，其中 $c
= 1$（取 $z = 1$）。记 $L = u + \iu v$，则 $\abs z = \eu^u$
且 $z = \abs z\eu^{\iu v}$，其中 $v$ 连续，$v(1) = 0$，
$v \in \intoo{-\pi}\pi$（$v$ 是 $z$
在连通的 $\Omega$ 上的一个连续辐角，于是其像避开奇数
倍的 $\pi$——$\Omega$ 中无点落在 $\R_-$ 上——
并含 $v(1) = 0$，故落在 $\intoo{-\pi}\pi$ 内：
$v$ 即主辐角），故 $L = \log$。

(b) 若 $g$ 是 $\C^*$ 上的一个连续对数函数：$h(t) =
g(\eu^{\iu t})$ 满足 $\eu^{h(t)} = \eu^{\iu t}$，故 $h(t)
- \iu t \in 2\iu\pi\Z$，由连续性得 $h(t) = \iu t +
2\iu\pi k$，其中 $k$ 为固定整数。于是 $g(1) = h(0) =
2\iu\pi k$，又 $g(1) = h(2\pi) = 2\iu\pi(k + 1)$：
矛盾。

(c) 在 $D(0,1)$ 上：$\log(1 + z) =
\sum_{n\geq1}\frac{(-1)^{n+1}}{n}z^n$——两边在
$0$ 处均为零，且导数都是 $\frac1{1+z} = \sum(-1)^nz^n$
（见\cref{ex:b3:holomorphic:examples}）；在连通圆盘上，
原函数在相差一个常数的意义下唯一。
\end{solution}

\begin{solution}{exo:b3:holomorphic:4}
(a) 在 $0$ 处展开（收敛半径为 $\infty$）：由\omterm{thm:b3:holomorphic:analytic}{柯西估计}，
在 $C_r$ 上 $\abs{c_k} \leq C(1 + r)^n/r^k \to 0$（当 $r \to
\infty$ 时，对 $k > n$）：$f = \sum_{k\leq n}c_kz^k$。

(b) 若 $\operatorname{Re}f \leq M$，则 $g = \eu^f$ 为整函数，且
$\abs g = \eu^{\operatorname{Re}f} \leq \eu^M$，由
刘维尔定理知其为常数。于是 $g' = f'g = 0$，且 $g$ 不取零值：$f' =
0$，故 $f$ 为常数。

(c) 若 $f$ 不取半平面 $H$ 中的值，则仿射映射 $w \mapsto
\alpha w + \beta$ 把 $\C\setminus H$ 映入
$\{\operatorname{Re} \leq M\}$ 中，对 $\alpha f +
\beta$ 应用 (b) 即可。
\end{solution}

\begin{solution}{exo:b3:holomorphic:5}
(a) $g(z) = f(z) - z^2$ 在点 $\frac1n$ 处为零，
这些点在 $0 \in \Omega$ 处聚集：由\omterm{thm:b3:holomorphic:identity}{恒等定理}
（$\Omega$ 连通），$g \equiv 0$，即 $f(z) = z^2$。

(b) 存在：$f(z) = z\cos(\pi/z)$ 在 $\C^*$ 上全纯
（复合运算），且 $f(\frac1n) = \frac1n\cos(n\pi) =
\frac{(-1)^n}n$。这与 (a) 并不矛盾：聚集
点 $0$ 是插值节点的聚集点，它\emph{不}属于
$\C^*$，故\omterm{thm:b3:holomorphic:identity}{恒等定理}在此不起作用——两个不同的
函数（例如 $z\cos(\pi/z)$，以及另一插值所得的
函数）可在这些点上取相同的值。

(c) $f \equiv 0$ 在全平面上，而 $g = 0$ 于 $D(0,1)$ 上，且
$g = 1$ 于 $D(3,1)$ 上：该函数在不连通的并集上全纯，
在 $D(0,1)$ 上相等，但二者并不相同。\omterm{thm:b3:holomorphic:identity}{恒等定理}的
开闭论证需要连通性，才能把结论由一个
分支传到另一分支——而这里做不到。
\end{solution}

\begin{solution}{exo:b3:holomorphic:6}
(a) 由有界区域上的最大模原理：
$\sup_{\mathbb D}\abs f = \sup_{\partial\mathbb D}\abs f = 1$。
因 $f$ 无零点，故 $1/f$ 在 $\mathbb D$ 上全纯，
在\omterm{def:b3:topology:interior}{闭包}上连续，且边界模为 $1$：同理
$\abs{1/f} \leq 1$，即 $\abs f \geq 1$。故 $\abs f \equiv
1$：模在\omterm{def:b3:topology:interior}{内部}取得最大值，且
由\cref{thm:b3:holomorphic:maximum}(2)可知 $f$ 为常数。

(b) $f(z) = z$：边界模为 $1$，在原点处有零点，
非常数——正是这一零点阻碍了 $1/f$
论证。
\end{solution}

\begin{solution}{exo:b3:holomorphic:7}
$G$ 的连续性：由控制收敛定理，控制函数为 $h$
（局部一致界）。全纯性由\omterm{thm:b3:holomorphic:weierstrassconv}{莫雷拉定理}推出（
在\cref{thm:b3:holomorphic:weierstrassconv}中建立）：对
圆盘内的闭三角形 $T$，其中 $\abs g \leq h$，
\[
\int_{\partial T}G(z)\,\dd z =
\int_X\Bigl(\int_{\partial T}g(x, z)\,\dd z\Bigr)\dd\mu(x) =
0,
\]
由富比尼定理交换积分次序（$\int_X\int_{\partial T}\abs
g \leq \operatorname{length}(\partial T)\int h < \infty$）及
内层积分由古尔萨引理为零。对于 $\Gamma$：在带形区域 $a
\leq \operatorname{Re}z \leq b$（$0 < a \leq b$），
$\abs{t^{z-1}\eu^{-t}} = t^{\operatorname{Re}z-1}\eu^{-t} \leq
(t^{a-1} + t^{b-1})\eu^{-t}$，在
$\intoo0{+\infty}$ 上\omterm{def:b3:lebesgue:l1}{可积}：$\Gamma$ 在
$\{\operatorname{Re} z > 0\}$ 上全纯（\omterm{def:b3:measure:measure}{测度}只是
$\sigma$-有限的，但证明只需\omterm{def:b3:lebesgue:l1}{可积}的
控制函数）。由\omterm{thm:b3:holomorphic:identity}{恒等定理}，函数方程
$\Gamma(z + 1) = z\Gamma(z)$，已在
$\intoo0{+\infty}$ 上证明（见\cref{ex:b3:lebesgue:gamma}），
在整个半平面上成立。
\end{solution}

\begin{solution}{exo:b3:holomorphic:8}
记 $P = c\prod_k(X - a_k)^{m_k}$
（见\cref{pb:b3:holomorphic:1}）。设 $P'(z) = 0$。若 $P(z) = 0$，
则 $z$ 为某个 $a_k$，位于凸包中。否则，
对数导数给出
\[
0 = \frac{P'(z)}{P(z)} = \sum_k\frac{m_k}{z - a_k}
= \sum_k m_k\,\frac{\bar z - \bar a_k}{\abs{z - a_k}^2} ;
\]
取共轭得 $\sum_kw_k(z - a_k) = 0$，其中 $w_k = m_k/\abs{z -
a_k}^2 > 0$：$z = \sum_k\frac{w_k}{\sum w}\,a_k$，即为根的凸
组合。对于 $P = z^3 - 1$，其根为
三次单位根，$P' = 3z^2$ 的二重根为 $0$——这正是
正三角形的重心。
\end{solution}

\begin{solution}{exo:b3:holomorphic:9}
闭单位正方形 $K = \{x + \iu y : 0 \leq x, y \leq
1\}$ 为紧集：$M = \sup_K\abs f < \infty$。对每个 $z \in \C$
与 $K$ 中某点的差属于 $\Z + \iu\Z$
（减去整数部分），且 $f$ 在这些平移下不变
（由两个关系式迭代可得），于是 $\abs f \leq M$ 在
$\C$ 上成立。由刘维尔定理，$f$ 为常数。因此任何非常值的
双周期亚纯函数——即经典理论中的椭圆
函数——必有极点。
\end{solution}

\begin{solution}{exo:b3:holomorphic:10}
(a) 在平均值公式中取实部
（见\cref{thm:b3:holomorphic:maximum}(1)）。光滑性：$f$ 是
解析函数，故 $P, Q \in \mathcal C^\infty$；对
\omterm{prop:b3:holomorphic:cauchyriemann}{柯西–黎曼方程}求导：$P_{xx} = (Q_y)_x = (Q_x)_y = (-P_y)_y =
-P_{yy}$（二阶偏导数的施瓦茨对称性）：$\Delta P =
0$。

(b) 若 $\operatorname{Re}f$ 在
连通的 $\Omega$ \omterm{def:b3:topology:interior}{内部}取得最大值：$g = \eu^f$ 有 $\abs g =
\eu^{\operatorname{Re}f}$ 在\omterm{def:b3:topology:interior}{内部}取得最大值，故
$g$，从而 $\operatorname{Re}f = \ln\abs g$，为常数
（见\cref{thm:b3:holomorphic:maximum}(2)）。在有界区域上，
并连续延拓至边界，则
$\sup_{\bar\Omega}\operatorname{Re}f =
\sup_{\partial\Omega}\operatorname{Re}f$。
\end{solution}

\begin{solution}{exo:b3:holomorphic:11}
(a) 两式在 $I$ 上相等（$z = \bar z$，且 $f$ 在该处取实值
：$\overline{f(\bar z)} = \overline{f(z)} = f(z)$），且
$z \mapsto \overline{f(\bar z)}$ 在
$\Omega^-\cup I$ 上连续（作为\omterm{def:b3:topology:continuity}{连续映射}的复合）：$F$ 在
$\Omega$ 上连续。$\Omega^-$ 上的全纯性：在 $z_0
\in \Omega^-$，将 $f(w) = \sum c_n(w - \bar z_0)^n$ 在
$\bar z_0 \in \Omega^+$ 附近展开，于是
\[
\overline{f(\bar z)} = \sum_n\bar c_n\,(z - z_0)^n,
\]
收敛幂级数，故全纯。

(b) 与 $I$ 不相交的三角形，由古尔萨引理在
$\Omega^\pm$ 中处理。对与 $I$ 相交的三角形，沿实轴将其
切分为至多三个三角形或四边形，每块各有一边
位于 $I$ 上；对这样的块 $P$，不妨设其包含于
$\overline{\Omega^+}$ 中，则\omterm{def:b3:holomorphic:contour}{围道积分}是
$\varepsilon \to 0^+$ 时，沿形如 $P + \iu
\varepsilon$ 型平移的积分的极限（$F$ 在紧块上一致连续，
使边界积分收敛，
$I$ 上的边自上方逼近），且每个平移
位于 $\Omega^+$ 中，古尔萨引理给出 $0$。对各
块求和，得 $\int_{\partial T}F = 0$。由\omterm{thm:b3:holomorphic:weierstrassconv}{莫雷拉定理}（准则
在\cref{thm:b3:holomorphic:weierstrassconv}中）：$F$ 在
$\Omega$ 上全纯。

(c) 在圆盘上，$G(z) = \overline{f(\bar z)}$
由 (a) 的计算知其全纯，且在直径上与 $f$
一致，直径是有聚点的集合，故 $G = f$ 处处成立
（由\omterm{thm:b3:holomorphic:identity}{恒等定理}）。若 $f$ 在非空\omterm{def:b3:topology:topology}{开集}
$U$ 上取实值，则在 $U$ 上 $Q = \operatorname{Im}f$ 的两个偏导数都
为零，\omterm{prop:b3:holomorphic:cauchyriemann}{柯西–黎曼方程}把这一结论转移到 $P =
\operatorname{Re}f$（$P_x = Q_y = 0$，$P_y = -Q_x = 0$），于是
$f' = P_x + \iu Q_x = 0$ 在 $U$ 上：$f$ 在 $U$ 上为常数，
由\omterm{thm:b3:holomorphic:identity}{恒等定理}，处处为常数（$\Omega$
连通）。
\end{solution}

\begin{solution}{exo:b3:holomorphic:12}
(a) $1 = f^2 + g^2 = (f + \iu g)(f - \iu g)$，故 $h = f +
\iu g$ 永不为零（否则其余因子须趋于无穷）。
对无零点的整函数 $h$：$h'/h$ 为整函数，且 $\C$ 是
星形的，故有原函数 $\varphi_0$
（见\cref{thm:b3:holomorphic:cauchy}）；于是
$\bigl(h\eu^{-\varphi_0}\bigr)' = \eu^{-\varphi_0}(h' -
h\varphi_0') = 0$：$h = c\,\eu^{\varphi_0}$，其中 $c \neq 0$，
并将常数 $\log c$ 吸收进 $\varphi = \varphi_0
+ \log c$（$c$ 的任一复对数）：$h =
\eu^{\varphi}$。

(b) 取 $h = \eu^{\varphi}$，且 $h^{-1} = f - \iu g =
\eu^{-\varphi}$：
\[
f = \frac{\eu^{\varphi} + \eu^{-\varphi}}2,
\qquad
g = \frac{\eu^{\varphi} - \eu^{-\varphi}}{2\iu} .
\]
记 $\varphi = \iu\psi$，其中 $\psi = -\iu\varphi$
为整函数，于是上式即 $f = \cos\psi$，$g = \sin\psi$。整函数解
恰为所有序对 $(\cos\psi,
\sin\psi)$，其中 $\psi$ 为整函数；逆命题则由
恒等式 $\cos^2 + \sin^2 = 1$ 给出。
对 $f^2 + g^2 = 0$：$(f + \iu g)(f - \iu g) = 0$，在
整环 $\mathcal H(\C)$ 中（\,$\C$ 连通：
零因子将违背\omterm{thm:b3:holomorphic:identity}{恒等定理}）：$g = \pm\iu
f$，其中 $f$ 为任意整函数。
\end{solution}

\begin{solution}{pb:b3:holomorphic:1}
\textbf{1.} 当 $\abs z \geq 1$ 时：
\[
\abs{P(z)} \geq \abs z^n\Bigl(1 -
\frac{\abs{a_{n-1}}}{\abs z} - \dots -
\frac{\abs{a_0}}{\abs z^n}\Bigr)
\geq \abs z^n\Bigl(1 - \frac{A}{\abs z}\Bigr),
\qquad A = \sum_k\abs{a_k} :
\]
当 $\abs z \geq R_0 = \max(1, 2A)$ 时，$\abs{P(z)} \geq
\frac12\abs z^n \to \infty$。

\textbf{2.} 取 $R \geq R_0$，使 $\frac12R^n \geq
\abs{P(0)}$。在紧集 $\bar D(0, R)$ 上，连续函数
$\abs P$ 于某点 $z_0$ 取得最小值；在圆外，$\abs P
\geq \frac12R^n \geq \abs{P(0)} \geq \abs{P(z_0)}$，因此
此最小值为全局的。

\textbf{3.} $Q(h) = P(z_0 + h)/P(z_0)$ 是 $h$ 的多项式，
且 $Q(0) = 1$；它非常值（$P$ 非常值），故某个
除常数项之外的系数不为零：$Q(h) = 1 +
c_kh^k + h^{k+1}S(h)$，其中 $k \geq 1$ 取最小，$c_k \neq 0$，
$S \in \C[X]$。

\textbf{4.} 关于根：任一 $w = \rho\eu^{\iu\varphi} \neq 0$ 有
第 $k$ 次方根 $\rho^{1/k}\eu^{\iu\varphi/k}$，其中
$\rho^{1/k}$ 存在，由介值定理用于
$t \mapsto t^k$ 在 $\intco0\infty$ 上——无循环论证。
取 $\omega$ 使 $\omega^k = -1/c_k$。于是
\[
Q(t\omega) = 1 - t^k + t^{k+1}\,\omega^{k+1}S(t\omega),
\qquad
\abs{Q(t\omega)} \leq 1 - t^k + C\,t^{k+1}
\quad (0 < t \leq 1),
\]
其中 $C = \abs\omega^{k+1}\sup_{\abs h \leq
\abs\omega}\abs{S(h)}$（注意 $1 - t^k \geq 0$ 在 $\intcc01$ 上成立）。

\textbf{5.} 对 $0 < t < \min(1, 1/C)$，有 $\abs{Q(t\omega)}
\leq 1 - t^k(1 - Ct) < 1$，即 $\abs{P(z_0 + t\omega)} <
\abs{P(z_0)}$——与全局极小性矛盾。故
$P(z_0) = 0$：达朗贝尔--阿尔冈证明至此完成。

\textbf{6.} 刘维尔版本：若 $P$ 无零点，则 $1/P$
为整函数；由第 1 题，$\abs{1/P} \leq 2R_0^{-n}$ 在
$\bar D(0, R_0)$ 之外，且 $1/P$ 在该紧圆盘上连续，
故在圆盘上亦有界：有界整函数为常数
（\cref{cor:b3:holomorphic:liouville}），从而 $P$ 为常数：
荒谬。要素：达朗贝尔证明用到紧性（最小值的存在）
及 $k$ 次根的\omterm{ex:b3:product:polar}{极坐标}形式存在性；刘维尔证明则用到
整函数的全套柯西理论（古萨定理——
本身亦是嵌套紧集论证——以及\omterm{thm:b3:holomorphic:analytic}{柯西
估计}），外加同样的强制性。闭圆盘的紧性
是两者共同的、\omterm{def:b3:representations:rep}{不可约}的核心。

\textbf{7.} 若 $\deg P \geq 1$，取一根 $\alpha$
（第 5 题）；作除法：$P = (X - \alpha)Q + P(\alpha) = (X -
\alpha)Q$，其中 $\deg Q = n - 1$；归纳即得。合并相同
因子：$P = c\prod_i(X - \alpha_i)^{m_i}$，$\sum m_i = n$。

\textbf{8.} 对实多项式 $P$：$P(\bar\alpha) =
\overline{P(\alpha)} = 0$，且重数相同
（对分解式取共轭）：非实根成对出现，
贡献 $(X - \alpha)(X - \bar\alpha) = X^2 -
2\operatorname{Re}(\alpha)X + \abs\alpha^2$，为一个实二次式，
判别式 $< 0$。由此得所述\omterm{def:b3:representations:rep}{不可约}多项式的
列表；又奇数次实多项式的非实根个数为偶数，
故必有一个实根。直接证法：
$P(x) \to \pm\infty$（当 $x \to \pm\infty$ 时，奇次，
首项系数设为正），故 $P$ 变号，并
于是介值定理适用。对 $X^3 - X - 1$：
两种论证都给出唯一实根 $\approx 1.3247$
（以及一对共轭复根）。

\textbf{9.} (i) $\chi_u \in \C[X]$ 为非常值多项式：它有一根
$\lambda$，且 $\det(u - \lambda\,\mathrm{id}) = 0$
给出一个特征向量；由
\cref{thm:b3:modules:jordan}中的\omterm{thm:b3:modules:structure}{初等因子}理论即适用于任一复矩阵，
其中 $\chi$ 总可分裂。(ii) 设 $P \in \bar\Q[X]$
为非常值多项式。作为 $\C$ 上的多项式，它有根 $z
\in \C$；$z$ 是 $\bar\Q$ 上的\omterm{def:b3:galois:algebraic}{代数元}，故由
传递性（\cref{cor:b3:galois:algclosed}）亦为 $\Q$ 上\omterm{def:b3:galois:algebraic}{代数元}，故 $z \in
\bar\Q$：$\bar\Q$ 上任一非常值多项式在
$\bar\Q$ 中有根。

\textbf{10.} 在 $\Q(\iu)$ 上，第 4 题已不成立：
$k$ 次根未必存在（无 $\sqrt2$），即使承认
根的存在性，则第 2 题亦不成立——极小化序列未必
收敛，因 $\Q$ 缺少完备性；在刘维尔的路线中，
古萨的嵌套紧三角形之交为空，
即在 $\Q(\iu)$ 的点上。两种证明都依赖 $\R$ 的完备性
（等价地，经有界单调收敛，即序完备性）；连通性
为\omterm{ex:b3:product:polar}{极坐标}形式背后的
介值定理提供动力。该命题
“$\C$ \omterm{def:b3:galois:closure}{代数闭}”是代数；已知的每一种
证明都是借 $\R$ 的定义偷渡进来的
分析。

\textbf{11.} 由\cref{thm:b3:holomorphic:analytic}，在 $a$ 处
$f^{(n)}(a) = n!\,c_n$，其中
\[
c_n = \frac1{2\iu\pi}\int_{C_r}\frac{f(w)}{(w -
a)^{n+1}}\,\dd w,
\qquad
\abs{c_n} \leq \frac{2\pi r}{2\pi}\cdot
\frac{\sup_{C_r}\abs f}{r^{n+1}}
= \frac{\sup_{C_r}\abs f}{r^{n}} :
\]
即\omterm{thm:b3:holomorphic:analytic}{柯西估计}（上式乘以 $n!$ 后的形式）。
若 $\abs f \leq M$ 在 $\C$ 上成立，则对任意
$a$ 及每个 $r$，有 $\abs{f'(a)} \leq M/r \to 0$，当 $r \to
\infty$ 时，故 $f' \equiv 0$，$f$ 在
连通的 $\C$ 上为常数——刘维尔定理由此重得。

\textbf{12.} 将 $f = \sum_kc_kz^k$ 在 $0$ 处展开（收敛半径
$\infty$）。对 $k > m$，有 $\abs{c_k} \leq (A + Br^m)/r^k \to
0$，当 $r \to \infty$ 时，故 $c_k = 0$，$f =
\sum_{k\leq m}c_kz^k$ 是次数至多为 $m$ 的多项式。

\textbf{13.} $g = \eu^f$ 为整函数，且 $\abs g =
\eu^{\operatorname{Re}f} \leq \eu^M$，由
第 11 题知其为常数。于是 $0 = g' = f'g$，又 $g$ 无零点，故 $f' = 0$，从而 $f$
为常数。

\textbf{14.} 设 $M = \sup_K\abs f$，$K$ 为紧的闭
单位正方形。每个 $z$ 与 $K$ 中某点之差
属于 $\Z + \iu\Z$ 中的元素（减去整数部分即可），并
反复使用两个周期性关系，
使 $f$ 保持不变，于是 $\abs f \leq M$ 在整个 $\C$ 上成立，由
第 11 题知 $f$ 为常数。因此，在格点下不变的非常值函数
不可能为整函数：经典理论中的椭圆函数
必须带有极点——这是通向
\cref{ch:b3:residues}的历史门径。

\textbf{15.} 若 $f(\C)$ 不与圆盘 $D(a, r)$ 相交，则
$\abs{f(z) - a} \geq r$ 对一切 $z$ 成立，故 $g = 1/(f - a)$ 为
整函数，且 $\abs g \leq 1/r$，由刘维尔定理知为常数，从而
$f$ 为常数。逆否命题：非常值整函数的值域
与每个圆盘都相交——即它在 $\C$ 中稠密。

\textbf{16.} 取 $R$，使 $\abs f \geq 1$ 在
$D(0, R)$ 之外成立。$f$ 的零点位于紧的
闭圆盘 $\bar D(0, R)$ 内；若它们有无穷多个，则必在其中聚集，
由\cref{thm:b3:holomorphic:identity}，将迫使 $f
\equiv 0$，不可能。记零点为 $z_1, \dots, z_p$，重数依次为
$m_1, \dots, m_p$，令 $M = \sum m_i$，并记
$\Pi(z) = \prod_i(z - z_i)^{m_i}$。从幂级数中逐个提出各零点，
得 $g = f/\Pi$ 为整且无零点的函数。
当 $\abs z \geq \max(R, 2\max_i\abs{z_i})$ 时：$\abs{z - z_i}
\leq 2\abs z$，又 $\abs f \geq 1$，故 $\abs{1/g} =
\abs\Pi/\abs f \leq 2^M\abs z^M$；在余下的紧
圆盘上，$1/g$ 连续，从而有界：$\abs{1/g} \leq A +
B\abs z^M$ 处处成立。由第 12 题知 $1/g$ 为
多项式；它无零点，故由第 7 题知它是
非零常数 $c$，从而 $f = \frac1c\Pi$ 为多项式。
反之，由第 1 题，每个非常值多项式都
是真映射。而 $\eu^z$ 确实不在此列：沿负实轴 $\R_-$
$\abs{\eu^z} = \eu^x \to 0$，而 $\abs z \to \infty$。

\textbf{17.} 对\cref{thm:b3:holomorphic:formula}作参数化，
在中心处取 $w = a + r\eu^{\iu\theta}$，$\dd w = \iu
r\eu^{\iu\theta}\dd\theta$,
\[
f(a) = \frac1{2\iu\pi}\int_{C_r}\frac{f(w)}{w - a}\,\dd w
= \frac1{2\pi}\int_0^{2\pi}
f\bigl(a + r\eu^{\iu\theta}\bigr)\,\dd\theta ;
\]
取实部即得 $u$ 的平均值性质。若
$u$ 在连通区域
$\Omega$ 的内点处取最大值：$\abs{\eu^f} = \eu^u$ 亦取\omterm{def:b3:topology:interior}{内部}最大值，
故 $\eu^f$ 为常数，由
\cref{thm:b3:holomorphic:maximum}(2)，且 $u =
\log\abs{\eu^f}$ 为常数。在直到边界连续的
有界区域上，$\sup_{\bar\Omega}u =
\sup_{\partial\Omega}u$，与 $\abs f$ 的情形完全一致。

\textbf{18.} \emph{情形 $\abs a > r$。}取 $R$，使 $r < R
< \abs a$：在凸圆盘 $D(0, R)$ 上，函数 $a - z$
全纯且无零点，而 $z \mapsto -1/(a - z)$ 在其上
存在原函数 $L$（见\cref{thm:b3:holomorphic:cauchy}）；
调整常数后，$\bigl(\eu^{-L}(a -
z)\bigr)' = \eu^{-L}\bigl(-L'\,(a - z) - 1\bigr) = 0$，得
$\eu^L = a - z$：存在全纯对数，且
$\log\abs{a - z} = \operatorname{Re}L(z)$。利用第 17 题的
平均值性质（取中心 $0$、半径 $r$）：
\[
\frac1{2\pi}\int_0^{2\pi}
\log\bigl|a - r\eu^{\iu\theta}\bigr|\,\dd\theta
= \operatorname{Re}L(0) = \log\abs a .
\]
\emph{情形 $\abs a < r$。}由 $a - r\eu^{\iu\theta} =
-r\eu^{\iu\theta}\bigl(1 - \frac ar\eu^{-\iu\theta}\bigr)$,
平均值等于 $\log r$ 加上 $\log\abs{1 -
\frac ar\eu^{-\iu\theta}}$ 的平均值。作代换 $\theta \mapsto
2\pi - \theta$，再——记 $\frac ar =
\rho\eu^{\iu\varphi}$，$\rho < 1$，$a = 0$ 的情形
平凡——平移 $\theta \mapsto \theta - \varphi$
（二者均保持一个周期上的平均值），把它化为
$\log\abs{1 - \rho\eu^{\iu\theta}}$ 的平均值：即第一种情形取
$(a, r) = (1, \rho)$，得 $\log 1 = 0$。综上，
两种情形下均有 $\log r = \log\max(\abs a, r)$。

\textbf{19.} $\log\abs{P(\eu^{\iu\theta})} = \log\abs c +
\sum_i\log\abs{\alpha_i - \eu^{\iu\theta}}$，每个根
按其重数重复计入；对 $\theta$ 取平均，并
应用第 18 题，取 $r = 1$，对每个不在单位
圆周上的根，得 $\log\bigl(\abs c\prod_i\max(1,
\abs{\alpha_i})\bigr)$。检验。对 $P = 2X - 1 = 2(X -
\frac12)$，则公式预测为 $\log 2$；直接计算，
$\abs{2\eu^{\iu\theta} - 1} = 2\abs{\frac12 -
\eu^{\iu\theta}}$，而 $\log\abs{\frac12 -
\eu^{\iu\theta}}$ 的平均值为 $\log\max(\frac12, 1) = 0$，故平均值为 $\log
2$。对 $P = X^2 - X = X(X - 1)$，根 $1$ 恰\emph{在}
圆周上；公式预测为 $0$。直接计算，
$\log\abs{\eu^{\iu\theta}}$ 的平均值为 $0$，并利用
$\abs{\eu^{\iu\theta} - 1} = 2\abs{\sin\frac\theta2}$:
\[
\frac1{2\pi}\int_0^{2\pi}
\log\Bigl(2\sin\frac\theta2\Bigr)\dd\theta
= \frac1\pi\int_0^\pi\log(2\sin u)\,\dd u
= \log 2 + \frac J\pi,
\qquad J = \int_0^\pi\log\sin u\,\dd u .
\]
作代换 $u = 2v$，并用 $\sin 2v = 2\sin v\cos v$，得
$J = \pi\log2 + 2\int_0^{\pi/2}\log\sin +
2\int_0^{\pi/2}\log\cos = \pi\log 2 + 2J$（每一半都等于
$J/2$，由 $\sin$ 的对称性），故 $J = -\pi\log 2$（
广义积分收敛，$\log\sin$ 在端点处\omterm{def:b3:lebesgue:l1}{可积}）：
即平均值为 $\log2 - \log2 = 0$。该公式
对位于圆周上的根依然成立。

\textbf{20.} $\frac{\eu^z - 1}z =
\sum_{j\geq0}\frac{z^j}{(j+1)!}$ 是整函数，且其值为 $1$，
在 $0$ 处：其倒数在 $0$ 附近全纯（在 $\abs z <
2\pi$ 内事实上成立，$\eu^z - 1$ 的最近其他零点为
$\pm2\iu\pi$），故 $\frac z{\eu^z-1}$ 在 $0$ 处解析。
将两个级数相乘，并读取
$z^n$ 的系数，$n \geq 1$，在 $(\frac{\eu^z-1}z)\cdot(\frac
z{\eu^z-1}) = 1$ 中，有：
\[
\sum_{k=0}^{n}\frac{B_k}{k!\,(n+1-k)!} = 0
\quad\Longleftrightarrow\quad
\sum_{k=0}^{n}\binom{n+1}{k}B_k = 0 .
\]
依次得 $B_0 = 1$，$B_1 = -\frac12$，$B_2 = \frac16$，
$B_3 = 0$，$B_4 = -\frac1{30}$，$B_5 = 0$，$B_6 =
\frac1{42}$。偶性：令 $F(z) = \frac z{\eu^z-1} +
\frac z2$，
\[
F(-z) = \frac{-z}{\eu^{-z} - 1} - \frac z2
= \frac{z\,\eu^z}{\eu^z - 1} - \frac z2
= z + \frac{z}{\eu^z - 1} - \frac z2 = F(z) :
\]
$F$ 为偶函数，故 $B_{2k+1} = 0$（$k \geq 1$，唯一的奇数指标
系数 $B_1$ 已被 $+\frac z2$ 吸收）。预告：
$\cot w = \iu + \frac{2\iu}{\eu^{2\iu w}-1}$ 可得
$w\cot w = 1 + \sum_{k\geq1}\frac{B_{2k}}{(2k)!}(2\iu
w)^{2k}$，于是余切函数的洛朗系数——
因此由\cref{ch:b3:residues}，每个 $\zeta(2k)$——
均由伯努利数确定：
\[
\zeta(2k) = (-1)^{k+1}\,\frac{(2\pi)^{2k}B_{2k}}{2\,(2k)!} .
\]

\textbf{21.} 设 $f = \sum_nc_nz^n$（收敛半径为 $\infty$）；
则 $g(z) = \overline{f(\bar z)} = \sum_n\bar c_nz^n$ 为
整函数。在 $\R$ 上，$g(x) = \overline{f(x)} = f(x)$，故 $g$ 与
$f$ 在连通的 $\C$ 中具有聚点的集合上相等：
由\cref{thm:b3:holomorphic:identity}得 $g \equiv f$，
即 $f(\bar z) = \overline{f(z)}$（等价地：所有
$c_n$ 为实数）。对实多项式 $P$：
$P(\bar\alpha) = \overline{P(\alpha)} = 0$，同样的
恒等式作用于实导数 $P', P'', \dots$
保持重数：非实根成对出现——
即练习 8 中的配对，此处以解析方法重新证明。

\textbf{22.} 将这一阶梯汇总如下：有界 $\Rightarrow$
常数（11）；受 $A + B\abs z^m$ 控制 $\Rightarrow$
多项式（12）；实部有上界 $\Rightarrow$
常数（13）；双周期 $\Rightarrow$ 常数（14）；
值域缺少一个圆盘 $\Rightarrow$ 常数（15）；真
映射 $\Rightarrow$ 多项式（16）。每一级都是柯西
公式：函数在一点处的值即圆周平均，故所有
泰勒系数都由 $f$ 在大
圆周上的大小决定，而增长上界可使系数
整体为零。这类条件对实的 $\mathcal
C^\infty$ 函数毫无约束：一个\omterm{def:b3:lp:mollifier}{凸起函数}
（见\cref{thm:b3:lp:regularization}）有界、紧
支、极端地非常数，其在支集之外任一点处的导数
全部为零，而函数在其附近并非处处为零。光滑性
使不同点处的导数毫无耦合；全纯性
则把每一阶导数串联到远处圆周上的单一积分。
一个局部条件配上一个全局“报信者”——
这正是整函数遵守定律与阶数的
缘由。

\textbf{23.} 展开
$\abs{P(\eu^{\iu\theta})}^2 = \sum_{k,l}a_k\bar a_l
\eu^{\iu(k-l)\theta}$，取平均则消去一切 $k \neq
l$ 的项，平均值为 $\sum_k\abs{a_k}^2 =: N$。先设
$P$ 在单位圆周上无零点，从而 $\theta \mapsto
\log\abs{P(\eu^{\iu\theta})}$ 连续。不等式 $\log
t \leq t - 1$ 作用于 $t = \abs P^2/N$，
取平均后得：
\[
\frac1{2\pi}\int_0^{2\pi}\log\abs{P(\eu^{\iu\theta})}^2
\dd\theta - \log N
\;\leq\; \frac1N\cdot N - 1 = 0,
\]
故 $\abs P$ 的几何平均至多为 $\sqrt N$；
练习 19 将该几何平均确定为
$\abs c\prod_i\max(1, \abs{\alpha_i})$，即兰道
不等式。圆周上的根：取 $r > 1$，使其
异于每个 $\abs{\alpha_i}$；多项式 $P(rX)$ 的根为
$\alpha_i/r$（不在单位圆周上），系数为 $a_kr^k$，
满足该不等式；两边关于 $r$ 连续，
令 $r \to 1^+$ 即得一般情形。对 $X^2 -
X$：根为 $0$ 与 $1$，故左边为 $1$，而
右边为 $\sqrt{1 + 1} = \sqrt2$：成立，且有余量。

\textbf{24.} 记 $\eu^z - 1 = z\,g(z)$，其中 $g(z) =
\sum_{k\geq0}\frac{z^k}{(k+1)!}$ 为整函数，$g(0) = 1$。由于
$\eu^z = 1$ 恰在 $2\pi\iu\Z$ 上成立，$g$ 在
$D(0, 2\pi)$ 内无零点（当 $0 < \abs z < 2\pi$ 时因 $\eu^z - 1
\neq 0$，在 $0$ 处由 $g(0) = 1$ 得证），故 $h = 1/g$ 在
$D(0,2\pi)$ 上全纯，其在 $0$ 处的泰勒级数——
按定义即 $\sum\frac{B_n}{n!}z^n$——在
整个圆盘上收敛：$\rho \geq 2\pi$。若 $\rho > 2\pi$，则和 $S$
将在 $D(0,\rho)$ 上全纯，且它与 $z
\mapsto z/(\eu^z - 1)$ 在 $0 < \abs z < 2\pi$ 上相等；二者
在连通\omterm{def:b3:topology:topology}{开集} $D(0,\rho) \setminus
2\pi\iu\Z$ 上均全纯，故由\omterm{thm:b3:holomorphic:identity}{恒等定理}它们在该处相等。但
当 $z \to 2\pi\iu$ 时，$\abs{z/(\eu^z - 1)} \to \infty$
（分子 $\to 2\pi$，分母 $\to 0$），而 $S$
在 $2\pi\iu$ 处连续：矛盾。因此 $\rho = 2\pi$
恰成立，由阿达马公式得
$\limsup_n\abs{B_n/n!}^{1/n} = \frac1{2\pi}$；奇数
项系数自 $n = 3$ 起均为零，故 $\limsup$
由偶数指标承担，这即所述公式
取 $n = 2k$。数值上，当 $k = 6$ 时：$(2\pi)^{12} \approx
3.7858\cdot10^9$，且 $2\cdot12! = 958\,003\,200$，故
$2\,(2k)!/(2\pi)^{2k} \approx 0.25305$，与
$\abs{B_{12}} = \frac{691}{2730} \approx 0.25311$ 相比，
比值 $1.00025$，在所示位数内恰为 $\zeta(12)$：
留数演算公式 $\zeta(2k) =
(-1)^{k+1}\frac{(2\pi)^{2k}B_{2k}}{2\,(2k)!}$（见
\cref{ch:b3:residues}）解释了因子 $2$ 以及
那微小的超出量。

\textbf{25.} 柯西界：若 $\abs z > 1 + M$，其中 $M =
\max_k\abs{q_k}$，则
\[
\Bigl|\sum_{k<n}q_kz^k\Bigr|
\leq M\,\frac{\abs z^n - 1}{\abs z - 1}
< \frac{M}{\abs z - 1}\,\abs z^n \leq \abs z^n,
\]
于是 $\abs{Q(z)} > 0$，所有根都落在 $\overline D(0, 1+M)$ 内。
$P_j$ 的系数收敛，因而有界，
即存在某个 $M$，使所有 $P_j$（以及 $P$）的根都落在
紧集 $K = \overline D(0, 1 + M)$ 中。设 $v_j \in K^n$ 为一个
向量，按重数列出 $P_j$ 的根
（第 7 题）。$(v_j)$ 的每个子列都有一个进一步的
子列收敛于某个 $(\beta_1, \dots, \beta_n)$；
$\prod_i(X - \alpha_i^{(j)})$ 的系数在符号意义下，
即为 $v_j$ 的初等对称函数——
连续——，故沿该子列，它们收敛于
$\prod_i(X - \beta_i)$ 的系数；而它们
依假设收敛于 $P$ 的系数，故 $\prod_i(X -
\beta_i) = P$：$(v_j)$ 的每个子列极限都是
$P$ 的根向量的一个置换。若匹配
距离 $\delta_j = \min_\sigma\max_i\,
\abs{\alpha_i^{(j)} - \alpha_{\sigma(i)}}$ 不趋于
$0$，则某子列将保持 $\delta_j \geq \varepsilon$
而其根向量收敛于
$P$ 的根的一个置换——，沿该子列便迫使 $\delta_j \to 0$：
矛盾。因此根的多重集收敛。关于精确性：
$P_\varepsilon = (X - 1)^2 + \varepsilon$ 的根为 $1 \pm
\iu\sqrt\varepsilon$：二重根移动的量为
$\sqrt\varepsilon$，例如移动约 $10^{-2}$，取 $\varepsilon =
10^{-4}$ 时。一般地，若 $\alpha$ 为 $m$ 重根，则
在 $\alpha$ 附近有 $\abs{P(z)} \asymp \abs{z -
\alpha}^m$，故大小为 $\varepsilon$ 的扰动使根的聚团
大约移动 $\varepsilon^{1/m}$：赫尔德
连续性，指数为 $\frac1m$，且不能更好——
这也解释了为何数值求解器在二重根附近
只能保留一半的有效数字。
\end{solution}
